\documentclass[11pt]{article}
\usepackage{preamble}
\setlength{\parskip}{4pt}

\begin{document}

\daybanner{AP Statistics \textperiodcentered\ Topic 1.3 (CED 1.3) \textperiodcentered\ B: 2026-09-11 \textperiodcentered\ E: 2026-09-14 \textperiodcentered\ TEACHER KEY}{A Percent Needs Its Base}

\sectionbanner{CED TETHER}

{\small
\begin{itemize}[leftmargin=1.5em,itemsep=2pt,topsep=2pt]
  \item \textbf{Skill 2.B} --- Construct numerical or graphical representations of distributions.
  \item \textbf{Skill 2.A} --- Describe data presented numerically or graphically.
  \item \textbf{Enduring Understanding UNC-1} --- Graphical representations and statistics allow us to identify and represent key features of data.
  \item \textbf{LO UNC-1.A} --- Represent categorical data using frequency or relative frequency tables. [Skill 2.B]
  \item \textbf{EK UNC-1.A.1} --- A frequency table gives the number of cases falling into each category. A relative frequency table gives the proportion of cases falling into each category.
  \item \textbf{LO UNC-1.B} --- Describe categorical data represented in frequency or relative tables. [Skill 2.A]
  \item \textbf{EK UNC-1.B.1} --- Percentages, relative frequencies, and rates all provide the same information as proportions.
  \item \textbf{EK UNC-1.B.2} --- Counts and relative frequencies of categorical data reveal information that can be used to justify claims about the data in context.
\end{itemize}
}
\textbf{Essential question:} When does a count answer the question, and when does a percent --- and what does each one hide?\par
\textbf{DOK-3 skill:} 4.B \textperiodcentered\ \textbf{FRQ pattern:} \texttt{frequency-vs-relative-frequency-which-table-the-argument-needs}\par
\textbf{DOK rationale:} Part (c) requires matching a stated question to the table that can answer it (counts vs. relative frequencies), citing the numbers that settle it, and reasoning about what each club's chosen table conceals from a reader; the choice must be defended against the other table rather than read off, and either table can be argued for.\par

\frameworkphaseheader{Do Now (first take) $\rightarrow$ Explore (video + follow-along) $\rightarrow$ Exit (finish + turn in)}{1 $\rightarrow$ 3}{45}{%
\begin{itemize}[leftmargin=1.5em,itemsep=2pt,topsep=2pt]
  \item Board slide up at the bell. Read the council's question aloud once; do not say which table you prefer.
  \item Video follow-along from minute 5 (frequency vs. relative-frequency tables).
  \item Board slide back up at minute 33. Circulate; every time a student says a percent, ask `of how many?'
  \item Collect at the door. Score part (c) only, E/P/I. Either table can earn E if the reason is tied to the question.
\end{itemize}
}{%
\begin{itemize}[leftmargin=1.5em,itemsep=2pt,topsep=2pt]
  \item One-sentence first take with one number, before any teaching.
  \item Video follow-along (build and read frequency and relative-frequency tables).
  \item Finish (a)--(c) with the rules callout in front of them; exit line; turn in.
\end{itemize}
}{%
\begin{itemize}[leftmargin=1.5em,itemsep=2pt,topsep=2pt]
  \item If the session room holds 30 chairs, which table tells the council how many chairs each club fills?
  \item 83\% of what? 60\% of what? Say the base out loud.
  \item Why did Robotics choose counts and Chess choose percents? What would each look like on the other table?
  \item Rewrite the council's question so that the \emph{other} table is the right one to answer it.
\end{itemize}
}{Solo teacher. Circulate ~5 pairs. Do not name a table yourself during the finish window; ask what the question is asking for --- people or share.}

\begin{callout}{\IconWarn\ WATCH FOR}{calloutred}
\begin{itemize}[leftmargin=1.5em,itemsep=2pt,topsep=2pt]
  \item `Percents are always fairer' stated as a rule, with no link to what the council actually asked.
  \item Citing 83\% without ever saying it is 15 of 18 people.
  \item Comparing 27 Yes to 15 Yes as if the clubs were the same size.
\end{itemize}
\end{callout}

\sectionbanner{SCORING PART (c) --- E / P / I}

\textbf{Expected elements}
\begin{itemize}[leftmargin=1.5em,itemsep=2pt,topsep=2pt]
  \item \textbf{table-tied-to-question} (required) --- Commits to one table and ties the choice to what the council's question asks (how many students vs. what share of a club); either table accepted with that reason
  \item \textbf{cites-settling-numbers} (required) --- Cites the specific numbers that settle it (27 vs. 15 Yes, or 83\% vs. 60\%)
  \item \textbf{counts-hide-base-sizes} (required) --- Says the Robotics frequency table hides that the clubs are different sizes (45 vs. 18), so 27 Yes is only 60\% / 18 said No
  \item \textbf{percents-hide-small-n} (required) --- Says the Chess percent table hides that 83\% is 15 people out of only 18 (a percent needs its base)
\end{itemize}

{\small\renewcommand{\arraystretch}{1.25}
\noindent\begin{tabular}{|p{0.5in}|p{5.9in}|}
\hline
\textbf{E} & Commits to one table with a reason tied to the council's question, cites the settling numbers, AND names both hides (club sizes behind the counts; the small base behind the percent). \\ \hline
\textbf{P} & Commits with numbers but names only one hide, OR names both hides but never commits to a table or gives no reason tied to the question, OR the reason is `percents are always fairer' with no link to what was asked. \\ \hline
\textbf{I} & Says the two tables say the same thing, or picks a table with no numbers, or names no hide. \\ \hline
\end{tabular}}

\textbf{Common mistakes}
\begin{itemize}[leftmargin=1.5em,itemsep=2pt,topsep=2pt]
  \item `Percents are always better' (or `counts are the real data') with no link to what the council asked
  \item Cites 83\% without ever saying it is 15 people out of 18
  \item Compares 27 Yes to 15 Yes without noting Robotics has 45 members to Chess's 18
  \item Reports 27/45 as `27\%' --- treats a count as a percent
  \item Rounds 15/18 to 80\% or 85\%
\end{itemize}

\bankitem{Optional \#1 --- not collected}{\dokbadge{2}~Suppose 25 students in a class were asked their favorite lunch: Pizza 11, Tacos 8, Salad 4, Other 2. Make the relative-frequency table (nearest whole percent) and check that your percents add to $100\%$. Then write one sentence about the class that a reader could verify from your table.}{}
\answer{Pizza $11/25 = 44\%$, Tacos $32\%$, Salad $16\%$, Other $8\%$; sum $= 100\%$. Sentence names the variable and a checkable claim, e.g. ``Fewer than half the class (44\%) chose pizza'' or ``Pizza and tacos together make up 76\% of the class.''}

\clearpage
\focusbanner{TODAY'S PROBLEM --- annotated}

Two clubs asked the student council for money to run a Saturday session. Each club surveyed \emph{all} of its members: ``Would you attend a Saturday session?'' The table shows every response as a count.\par\smallskip At the council meeting, the \textbf{Robotics club} presented a \textbf{frequency table} (counts) and the \textbf{Chess club} presented a \textbf{relative-frequency table} (percents). Each club chose the table that made its case look best.\par
\begin{center}
\textbf{``Would you attend a Saturday session?'' (every member surveyed; counts)}\par\smallskip
\begin{tabular}{|l|c|c|c|}
\hline
\textbf{Club} & \textbf{Yes} & \textbf{No} & \textbf{Total} \\ \hline
\textbf{Chess} & 15 & 3 & 18 \\ \hline
\textbf{Robotics} & 27 & 18 & 45 \\ \hline
\end{tabular}
\end{center}
\begin{firsttakebox}{FIRST TAKE (5 min)}
Which club has more interest in a Saturday session --- Chess or Robotics? One sentence and one number.\par
\answer{Not graded. Expect a split: some say Robotics (27) and some say Chess (83\%); the split IS the lesson.}
\end{firsttakebox}

\textit{Rules callout ``TABLES FOR A CATEGORICAL VARIABLE (keep on the page)'' is printed on the student sheet.}\par\medskip

\dokbadge{1}~\textbf{(a)} How many Robotics members said Yes? Which club has more members, and how many more?\par
\answer{27 Robotics members said Yes. Robotics has more members: 45 vs.\ 18, so 27 more.}

\dokbadge{2}~\textbf{(b)} Build the relative-frequency table: for each club, the percent of its members who said Yes and the percent who said No (nearest whole percent). Which club has the higher percentage saying Yes?\par
\answer{Chess: Yes $15/18 \approx 83\%$, No $3/18 \approx 17\%$. Robotics: Yes $27/45 = 60\%$, No $18/45 = 40\%$. Each row adds to $100\%$. Chess has the higher percentage saying Yes ($83\%$ vs.\ $60\%$).}

\dokbadge{3}~\textbf{(c)} The council asks: ``Which club has \textbf{more student interest} in a Saturday session?'' Decide which table --- counts or percents --- answers \emph{that} question, and cite the specific numbers that settle it. Then explain what each club's chosen table \textbf{hides} from a reader: what does the Robotics club's frequency table hide, and what does the Chess club's percent table hide?\par
\answer{Expected: ``more student interest'' asks \emph{how many students} are interested, so the \textbf{frequency table} answers it: 27 Robotics members said Yes vs.\ 15 Chess members --- 12 more people would show up. (A student who reads the question as ``which club's \emph{members} are more interested'' and picks the percent table, $83\%$ vs.\ $60\%$, earns full credit \emph{if they say that is the question the percents answer}.) What the tables hide: the Robotics club's \textbf{counts} hide the base sizes --- Robotics has 45 members to Chess's 18, so 27 Yes is only $60\%$ of the club and 18 Robotics members said No. The Chess club's \textbf{percents} hide how small the base is --- $83\%$ is 15 people out of only 18; $83\%$ of 18 (15) is fewer students than $60\%$ of 45 (27).}

\begin{summaryexitbox}{EXIT}
One thing the video changed about my first take: \hrulefill
\end{summaryexitbox}

\end{document}
